% ЛЕТНИЙ ОТЧЁТ (черновик). Стиль выдержан под зимний main.tex: простыми словами, не перегружено.
% Рисунки кладутся в подпапку figures/ (копируются из analysis/metrics/figures/).
% Заглушки [В РАБОТЕ] помечают результаты, которые ещё досчитываются (data-scaling, чистый T3).
\documentclass{article}
\usepackage[utf8]{inputenc}
\usepackage[T2A]{fontenc}
\usepackage[russian]{babel}
\usepackage{graphicx}
\usepackage{amsmath}
\usepackage{url}
\usepackage{geometry}
\geometry{a4paper, margin=2cm}

% мягкая заглушка для ещё не готовых рисунков
\newcommand{\MaybeImage}[2]{%
  \IfFileExists{#1}{\begin{center}\includegraphics[width=#2\linewidth]{#1}\end{center}}%
  {\begin{center}\fbox{\parbox[c][2.2cm][c]{0.8\linewidth}{\centering\small Рисунок в работе:\\\texttt{\detokenize{#1}}}}\end{center}}%
}

\title{Прогнозирование расщепления белков на пептиды}
\author{Шабаров Игорь}
\date{Июнь 2026}

\begin{document}
\maketitle

\section{Актуальность и значимость}
\textbf{Белки} — это биополимеры из аминокислотных остатков; они выполняют в клетке структурные, каталитические и регуляторные функции. Многие белки синтезируются в виде неактивного предшественника, который затем разрезается специальными ферментами — \textbf{протеазами}. В ходе этого расщепления (\textbf{протеолиза}) образуются \textbf{пептиды} — короткие фрагменты, как правило до нескольких десятков остатков, и нередко именно они несут биологическую активность: участвуют в сигнальных и иммунных процессах, регуляции клетки.

В геномных и протеомных базах обычно записана последовательность белка-предшественника, а не конечные продукты его разрезания. Поэтому задача предсказать, на какие пептиды распадётся предшественник, важна сама по себе: она позволяет восстановить набор потенциально образующихся пептидов и точнее интерпретировать данные.

Где это применяется:
\begin{itemize}
  \item \textbf{Протеомика.} Сопоставление наблюдаемых масс-спектрометрией пептидов с белками-предшественниками и локализация мест разреза.
  \item \textbf{Терапевтические пептиды и вакцины.} Оценка «обрабатываемости» конструкции: какие фрагменты вероятно образуются \textit{in vivo}, какие паттерны разреза ведут к нежелательным продуктам.
  \item \textbf{Дизайн белков.} Предотвращение нежелательного расщепления при проектировании последовательностей.
  \item \textbf{Патологии.} В ряде заболеваний протеолиз нарушается; модели помогают сопоставлять ожидаемые и наблюдаемые паттерны разреза.
\end{itemize}

\section{Постановка задачи}
Дан белок-предшественник — последовательность аминокислот
$$ x = (a_1, a_2, \dots, a_L), \qquad a_t \in \Omega, $$
где $\Omega$ — алфавит из 20 аминокислот, $L$ — длина. Нужно предсказать, какие непрерывные участки этой последовательности станут после разрезания зрелыми \textbf{пептидами}, а какие — \textbf{пропептидами} (вспомогательными участками, которые отрезаются и удаляются, но сами не являются конечным активным продуктом).

Удобно сформулировать это как разметку каждой позиции одной из трёх меток
$$ y_t \in \{\textsf{None},\ \textsf{Peptide},\ \textsf{Propeptide}\}. $$
Непрерывные участки одинаковой метки задают \textbf{сегменты} — пары (начало, конец) с типом. То есть ответ модели — это набор типизированных отрезков на последовательности. Например:
$$ \texttt{MGKRALRRSGKVRNL} \;\rightarrow\; \texttt{MGK} \mid \texttt{RALR} \mid \texttt{RSGK} \mid \texttt{VR} \mid \texttt{NL} $$

Что считается \textbf{правильным предсказанием.} Предсказанный пептид засчитывается как верный (true positive), если \emph{и} его начало, \emph{и} его конец совпадают с началом и концом настоящего пептида того же типа с точностью до небольшого допуска $\pm\tau$ (мы используем $\tau=3$ остатка). Предсказанные сегменты без пары — ложные срабатывания, настоящие сегменты без пары — пропуски. Качество измеряем по точности (precision), полноте (recall) и их среднему гармоническому (F1) отдельно для пептидов и пропептидов. Допуск $\pm\tau$ отражает то, что биологически важно угадать участок разреза, а не каждую позицию идеально.

\section{Обзор существующих решений}
Для этой задачи выделим два класса методов: специализированные модели предсказания протеолитических пептидов и общие \textbf{белковые языковые модели} (protein language models, pLM), которые дают признаки последовательности.

\subsection{Специализированные модели}
Одна из немногих моделей, прямо нацеленных на задачу, — \textbf{DeepPeptide}~\cite{deeppeptide2023}. Она берёт аннотации типов \texttt{PEPTIDE} и \texttt{PROPEP} из \textbf{UniProtKB/Swiss-Prot}~\cite{uniprot2025} и решает задачу как сегментацию предшественника. Её устройство:
\begin{itemize}
  \item эмбеддинги из языковой модели ESM-2~\cite{esm2_2019};
  \item стек CNN--BiLSTM--CNN поверх по-остаточных эмбеддингов;
  \item линейно-цепной CRF с расширенным набором состояний, который собирает целостные сегменты с ограничением на длину.
\end{itemize}
Такой декодер оптимизирует качество на уровне целых пептидов, а не отдельных позиций. Мы берём DeepPeptide как сильную отправную точку и улучшаем её.

\subsection{Языковые модели как источник эмбеддингов}
\begin{itemize}
  \item \textbf{ESM-2}~\cite{esm2_2019} — трансформер-энкодер, обученный предсказывать замаскированные остатки; даёт контекстные по-остаточные эмбеддинги.
  \item \textbf{ESM Cambrian (ESM C)}~\cite{esm_cambrian_blog_2024} — более новое семейство энкодеров, оптимизированное по качеству и ресурсам; рассматривается как замена ESM-2. Старшая версия — на 6 миллиардов параметров.
  \item \textbf{Ankh}~\cite{ankh2023} — компактная encoder--decoder модель, на практике используют только энкодер.
\end{itemize}
pLM применяются как «замороженные» энкодеры: последовательность переводится в матрицу эмбеддингов $L \times d$, поверх которой обучается модель сегментации.

\subsection{Структурные и генеративные модели}
Структурные модели, в частности \textbf{AlphaFold~3}~\cite{alphafold3_2024}, восстанавливают трёхмерную структуру и дают пространственный контекст, недоступный из одной последовательности. Генеративные модели (\textbf{DiMA}~\cite{dima2024}) ориентированы на дизайн новых пептидов и могут служить на этапе генерации кандидатов. В нашей работе мы используем структурный сигнал косвенно — через предсказанный «структурный алфавит» 3Di (см.\ ниже).

\subsection{Гомологичное разбиение и оценка}
Чтобы оценка не была завышена за счёт похожих последовательностей в обучении и тесте, данные разбивают на фолды так, чтобы между фолдами не было гомологии. Для этого применяют \textbf{GraphPart}~\cite{graphpart2023}: он разбивает последовательности на фолды с ограничением на попарную идентичность (мы используем порог 30\%). Оригинальный DeepPeptide оценивался вложенной перекрёстной проверкой (nested cross-validation) на 5 фолдах.

\section{Базовая модель и её недочёты}
DeepPeptide — разумная архитектура, но не идеальная. В ходе работы мы нашли несколько моментов, которые важно учесть.

\paragraph{Метрика разреза.} В исходном коде подсчёта метрик ($\pm3$) обнаружилась ошибка с перекрытием имён переменных, из-за которой полнота систематически занижалась. Мы её исправили; во всём отчёте используем исправленный подсчёт. Поправка сдвигает числа примерно на $+0.01$--$0.015$ почти равномерно и не меняет порядок моделей, то есть на выводы не влияла.

\paragraph{Как выбиралась архитектура.} Вложенная перекрёстная проверка честно настраивает гиперпараметры, не заглядывая в тестовый фолд. Но сам \emph{выбор архитектуры} в оригинале делался на тех же фолдах, на которых потом приводятся результаты, а итоговая модель для применения — ансамбль из 20 сетей — на честно отложенных данных не измерялась. Поэтому опубликованное число — слегка оптимистичная оценка. Это наблюдение и подтолкнуло нас построить более аккуратный протокол оценки (раздел~\ref{sec:protocol}).

\section{Данные: от 2022 к 2026}
Исходный датасет DeepPeptide собран из Swiss-Prot (версия 2022 года, $\approx$7{,}6 тыс.\ белков). Мы пересобрали датасет по той же рецептуре на свежей версии (2026 год): белков стало заметно больше, добавились новые семейства. Данные сильно неоднородны по организмам (от человека и мыши до ядовитых конусов и пауков) и по длинам пептидов (от 5 до 50 остатков). Эта неоднородность — важный сюжет всего отчёта: разные семейства ведут себя по-разному, и это влияет на сравнение моделей (раздел~\ref{sec:protocol}).

\section{Чистый протокол оценки}\label{sec:protocol}
Это — основная методическая часть работы.

\paragraph{Разбиение и роли фолдов.} Мы воспроизвели рецептуру статьи: разбили данные 2026 года на 7 фолдов с помощью GraphPart (порог идентичности 30\%) и сбалансировали фолды по флангующим мотивам мест разреза. Затем явно закрепили \emph{роли}:
\begin{itemize}
  \item \textbf{train} (4 фолда) — обучение;
  \item \textbf{dev} (1 фолд) — выбор эпохи (ранняя остановка);
  \item \textbf{model-select} (1 фолд) — сравнение архитектур;
  \item \textbf{sealed} (1 фолд) — изначально планировался как запечатанный финальный тест.
\end{itemize}
Идея — отделить выбор архитектуры от итоговой оценки, чего не хватало оригиналу.

\paragraph{Проблема: фолды неэквивалентны.}\label{par:problem} Балансировка делалась по мотивам разреза, но \emph{не} по длинам пептидов и не по организмам. В итоге фолды заметно различаются по составу, и результаты на разных фолдах сильно расходятся: модели «меняются местами» в рейтинге от фолда к фолду, а отдельные семейства (например, яд паука \textit{Cyriopagopus}) на «своём» фолде проседают на десятки пунктов. Это эффект \emph{состава фолда}, а не самих моделей.

\MaybeImage{figures/fold_divergence.png}{0.9}

\paragraph{Полное решение и компромисс.} Полностью эту проблему лечит полная перекрёстная проверка (nested cross-validation): тогда каждый белок ровно один раз попадает в тест, и редкие семейства взвешены по своей настоящей частоте. Но это многократно дороже по вычислениям, и мы её пока не проводим. Вместо этого мы \textbf{объединяем (pool) два отложенных фолда} (model-select и sealed) и сравниваем модели на объединении. Удачно, что эти два фолда комплементарны по длине (в одном больше коротких пептидов, в другом — длинных), поэтому их перекосы взаимно гасятся, а размер выборки удваивается ($\approx$2{,}3 тыс.\ белков). Платой за это становится отсутствие фолда, «запечатанного» от выбора архитектуры, — но его не было и в оригинале DeepPeptide, так что мы честно отмечаем это как ограничение. Там, где важно показать, \emph{где} какая модель сильна, мы используем разбор по стратам (длина, тип, организм): доля верно найденных сегментов внутри страты не зависит от состава фолда, поэтому объединение здесь корректно.

\section{Эксперименты и результаты}
Все числа ниже — на объединённых отложенных фолдах, исправленная метрика $\pm3$. Сравнение моделей показано на рис.~\ref{fig:scoreboard}.

\begin{figure}[ht]\centering
\MaybeImage{figures/scoreboard.png}{1.0}
\caption{Сравнение моделей на объединённых отложенных фолдах: F1, точность, полнота (с 95\%-ми доверительными интервалами). Лесенка от ESM-2 до полной архитектуры на ESM-C 6B.}
\label{fig:scoreboard}
\end{figure}

\subsection{Данные — главный рычаг}
Самый сильный практический вывод: сильнее всего качество поднимают не архитектурные надстройки, а \textbf{объём и покрытие обучающих данных}. Последовательности, непохожие на обучающие (а при гомологичном разбиении отложенные фолды именно такие), предсказываются плохо; значит, чем больше данных, тем шире покрытие пространства последовательностей и выше качество. Мы воспроизводим это на чистом протоколе: обучаем на возрастающих долях обучающих данных и смотрим на кривую качества.

\MaybeImage{figures/datascale_curve.png}{0.8}

\noindent\textit{[В РАБОТЕ: кривая F1 от объёма обучающих данных для baseline и сильной модели; ожидаем, что кривая ещё растёт к 100\,\% — данные не исчерпаны.]}

\subsection{Архитектура: голова работает только на богатых эмбеддингах}
Мы добавляли к базовой архитектуре «голову границ» (boundary head), которая отдельно учится предсказывать места разреза. Три наблюдения вместе дают аккуратный вывод:
\begin{itemize}
  \item \textbf{Голова на богатых эмбеддингах (ESM-C 6B)} даёт устойчивый прирост $+0.05$--$0.07$ F1 на обоих фолдах.
  \item \textbf{Богатые эмбеддинги сами по себе} почти не помогают: «голый» ESM-C 6B обгоняет ESM-2 лишь на $\approx$0.03, и этот выигрыш неустойчив (на разных фолдах меняет знак).
  \item \textbf{Голова на слабых эмбеддингах (ESM-2)} не даёт почти ничего ($\approx 0$). \textit{[В РАБОТЕ: чистый замер без bond-лосса.]}
\end{itemize}
Вывод: рычаг — это \emph{взаимодействие} «голова $\times$ богатые эмбеддинги», ни то, ни другое по отдельности. Это и объясняет, почему наши модификации «раскрылись» на ESM-C, но почти ничего не дали на ESM-2.

\MaybeImage{figures/interaction.png}{0.9}

\paragraph{Как и где помогает голова.} Голова заостряет границы: если сжимать допуск разреза с $\pm3$ до точного совпадения, модели с головой теряют качество гораздо медленнее (доля сохранённого F1 растёт по лесенке примерно с $0.55$ у базовой модели до $0.72$). Прирост головы сосредоточен на полноте у длинных пептидов и на точности у коротких предсказаний.

\subsection{Дополнительные каналы — это размены, а не рычаги}
\begin{itemize}
  \item \textbf{Структурный канал 3Di} (предсказанный структурный алфавит) — это \emph{размен}: помогает пропептидам, но вредит пептидам (в основном по полноте, особенно на длинных). В сумме на объединённых фолдах — примерно ничья.
  \item \textbf{Сжатие эмбеддинга} с 2560 до 256 признаков практически не меняет качество (разница в пределах шума). Это значит, что модель можно сделать в $\approx$16 раз компактнее без потерь — то есть выбор сжатия обоснован, а не произволен.
  \item \textbf{Bond-лосс} (дополнительный штраф за границы) на чистых данных оказался скорее нейтрально-вредным: в основном теряет полноту на пептидах. Это не чистый размен, а близко к бесполезному.
\end{itemize}
Сквозной мотив: и 3Di, и bond имеют одну «подпись» — сдвигают модель в сторону пропептидов ценой полноты на пептидах.

\MaybeImage{figures/trades.png}{0.9}

\subsection{Прочие проведённые эксперименты}
Мы честно перепроверили большой набор более ранних идей (есть сводная таблица с исправленной метрикой), но они оценивались по \emph{старому} методу — на разбиении 2022 года и с выбором модели по валидации, — поэтому требуют перепроверки на чистом протоколе. Среди них: ширина структурной проекции (16/32/48 — без эффекта), телескопический сегментный CRF (незначимо), варианты слияния со словарём пептидов (Aho), структурные признаки из AlphaFold, дообучение ESM-2 через LoRA, мультимасштабные проекторы. Мы не замалчиваем эту работу, но и не делаем по ней выводов без повторной проверки.

\section{Выводы}
\begin{itemize}
  \item Главный практический рычаг — \textbf{данные}: качество ограничено их объёмом и покрытием, непохожие на обучение последовательности предсказываются плохо.
  \item Единственная \emph{архитектурная} надстройка с устойчивым выигрышем (из проверенных) — \textbf{голова границ на богатых эмбеддингах}; сам по себе ни более сильный эмбеддинг, ни голова на слабом эмбеддинге устойчивого прироста не дают.
  \item Структурный канал 3Di и bond — это познавательные размены, а не безусловная польза; сжатие эмбеддинга безвредно и позволяет уменьшить модель.
  \item Единого «чемпиона» поверх этого нет: на фолдах разного состава лидер меняется, поэтому мы говорим аккуратно — «здесь сильнее эта модель, там другая».
\end{itemize}

\section{Ограничения и дальнейшая работа}
\begin{itemize}
  \item Нет полной перекрёстной проверки — это единственный способ получить одно чистое глобальное число с честным выбором архитектуры; отложено из-за стоимости.
  \item Нет фолда, запечатанного от выбора архитектуры (следствие объединения фолдов).
  \item Доверительные интервалы учитывают только разброс по белкам (не по сидам/обучению), поэтому оценка консервативна.
  \item Заявление об обобщении на новые семейства мы делаем узко (полнота не падает на белках, добавленных после 2022 года, внутри отложенного набора); честный кросс-датасетный тест (обучение на 2022, проверка на 2026) — в планах.
  \item Не проверены на чистом протоколе: Ankh, дообучение pLM и контрастное обучение, трансформер-голова поверх эмбеддингов.
\end{itemize}

\bibliographystyle{plain}
\bibliography{refs}

\end{document}
